\documentclass[10pt,a4paper]{amsart}
\usepackage[margin=19mm]{geometry}
\usepackage{amsmath,amssymb,mathtools}
\usepackage[scr=rsfs,cal=euler]{mathalfa}
\usepackage{xfrac}
\usepackage[unicode,hidelinks,bookmarks=false]{hyperref}
\newcommand{\ie}{i.e.\ }
\newcommand{\cf}{cf.\ }
\newcommand{\resp}{respectively\ }
\newcommand{\Subsection}{\subsection}
\providecommand{\nobreakdash}{\nobreak\hbox{-}}
\setlength{\emergencystretch}{2em}
\multlinegap0pt
\usepackage{xcolor}
\usepackage{enumerate}
\usepackage[shortlabels]{enumitem}
\usepackage{dsfont}
\usepackage{amssymb}
\usepackage{mathtools}
\newtheorem{lemma}{Lemma}
\title{Locally finite Fixed points of Branching Brownian Motion}
\author{Xinxin Chen, Arnab Chowdhury, Atul Shekhar, Shuo Zhu}
\date{Source: arXiv:2608.23202v1 (CC BY 4.0)}
\begin{document}
\maketitle

\noindent\emph{Source and licence.} Locally finite Fixed points of Branching Brownian Motion. Version arXiv:2608.23202v1 (CC BY 4.0), distributed by arXiv under the Creative Commons Attribution 4.0 International licence, http://creativecommons.org/licenses/by/4.0/, as recorded in the arXiv licence metadata for this exact version and shown on the abstract page https://arxiv.org/abs/2608.23202v1. Full licence notice: \texttt{licenses/arxiv-2608.23202-CC-BY-4.0.txt}.\par\smallskip

\noindent\textit{Editorial scope.} Counted unit: the article's lemma sec3 (source0.tex lines 1623--1628). The article's complete original proof of it is delivered with it (source0.tex lines 1629--1684, one \texttt{proof} environment of 56 lines). No mathematics is written, repaired or re-derived: the statement and the proof are the article's own lines, in the article's own notation, and no retained line is substituted or re-flowed.  Retained uncounted as the source-local context the proof needs: lines 476--482.  Nothing was omitted from any retained span, so the package carries no omission record.  No retained line contains a citation, so the wrapper carries no bibliography and no bibliography entry.  No retained line contains a figure environment, an image-inclusion command or drawing code, so no image is delivered and the package declares no asset dependency.  Editorial formatting only, in addition: an AMS \texttt{amsart} preamble that restates the article's own declarations and macros used by the retained lines, namely the environments \texttt{lemma} on separate counters, together with the standard packages those lines need.  The retained environments print in this document as Lemma 1 in order of appearance, whereas the article prints the same items with the numbering of its own full document.  The complete native source of the article as distributed in the arXiv e-print for this exact version is bundled unchanged as \texttt{sources/208397/source0.tex}.
\begin{lemma}[The many-to-one lemma]\label{Lemma many to one}
Let $F: \mathbb R \to \mathbb R$ be a bounded measurable function. Then, 
    \begin{equation}
        \mathbb E[\sum_{k \le n(t)} F(\chi_k(t))]=e^t \mathbb E[F(B_t)],
    \end{equation}
    where $B$ is a standard one dimensional Brownian motion.
\end{lemma}

\begin{lemma}\label{Main_contribution in sec3}
  Let $(p_{i})_{i\in I}$ be the atoms of a Poisson point process $\mathcal P$ with intensity measure $e^{-\mu x}dx$. For any choice of $z\in \mathbb R$ and $\epsilon>0$, there exist $t_{0}>0, \delta >0$ such that for all $t > t_{0}$, 
  \begin{equation}
    \mathbb{P} \biggl(\exists i\in \mathbb N, k \le n^{i}(t):p_{i}+\chi_k^i(t)-\lambda t \ge z,p_i\notin[-(\mu - \lambda)t-\delta\sqrt{t},-(\mu -\lambda)t+\delta\sqrt{t}]\biggr)\le \epsilon. \label{2.18}
  \end{equation}
\end{lemma}

\begin{proof}
  We first consider the case when $p_{i}> -(\mu -\lambda)t+\delta\sqrt{t}$.
  By union bounds and Campbell's formula, we have 
  \begin{equation}\label{estimae in Lem3.2}
    \begin{aligned}
       &\mathbb{P} \biggl(\exists i\in \mathbb N, k \le n^{i}(t):p_{i}+\chi_k^i(t)-\lambda t \ge z, p_i>-(\mu -\lambda)t+\delta\sqrt{t} \biggr)\\
     \le&\mathbb{E} \biggl[\sum_{p_{i} > -(\mu - \lambda)t +\delta \sqrt{t}}\mathbb{P} (p_{i}+M_{t}-\lambda t \ge z)\biggr]\\
      =&\int_{-(\mu-\lambda)t+\delta\sqrt{t}}^{\infty}\mathbb{P} (M_{t}\ge \lambda t-x+z)e^{-\mu x}dx.\\
    \end{aligned}
  \end{equation}
We use the many to one lemma  (Lemma \ref{Lemma many to one}) to bound the probability inside the integral and split it into two parts:
\begin{equation}
    \begin{aligned}
    &\int_{-(\mu-\lambda)t+\delta\sqrt{t}}^{\infty}\mathbb{P} (M_{t}\ge \lambda t-x+z)e^{-\mu x}dx\\
           \le&\int_{-(\mu-\lambda)t+\delta\sqrt{t}}^{\infty}e^{t}\mathbb{P}(B_{t}\ge \lambda t-x+z)e^{-\mu x}dx \\
    = &\int_{-(\mu-\lambda)t+\delta\sqrt{t}}^{0}e^{t}\mathbb{P}(B_{t}\ge \lambda t-x+z)e^{-\mu x}dx + \int_{0}^{\infty}e^{t}\mathbb{P}(B_{t}\ge \lambda t-x+z)e^{-\mu x}dx\\
     = &(I) + (II). \label{2.19}
    \end{aligned}
\end{equation}
By standard Gaussian Mills ratio argument, we get
  \begin{equation}
    \begin{aligned}
      (I)
     =&\int_{-(\mu-\lambda)t+\delta\sqrt{t}}^{0}e^{t}\mathbb{P}(B_{t}\ge \lambda t-x+z)e^{-\mu x}dx\\
\lesssim & \int_{0}^{(\mu-\lambda)t-\delta\sqrt{t}}\frac{e^{t}}{\sqrt{2\pi t}}e^{\mu x}e^{-\frac{(z+x+\lambda t)^{2}}{2t}}dx.\\
    \end{aligned}
  \end{equation}
  where we do a change of variables $x\to -x$. 
  Observe that since $\mu = \lambda + \sqrt{\lambda^2-2}$, we can find a constant $c_{1}>0$ depending on $z$ such that for  $t$ large enough,
  \begin{equation}
    \begin{aligned}
      &\int_{0}^{(\mu-\lambda)t-\delta\sqrt{t}}\frac{e^{t}}{\sqrt{2\pi t}}e^{\mu x}e^{-\frac{(z+x+\lambda t)^{2}}{2t}}dx\\
    \le& c_{1}\int_{0}^{(\mu-\lambda)t-\delta\sqrt{t}}\frac{1}{\sqrt{2\pi t}}e^{-\frac{(x-(\mu -\lambda)t)^{2}}{2t}}dx\\
     = & c_{1}\int_{-(\mu-\lambda)\sqrt{t}}^{-\delta}\frac{1}{\sqrt{2\pi}}e^{-\frac{u^{2}}{2}}du.\\
    \end{aligned}
  \end{equation}
  As for the second term, direct calculation gives that 
  \begin{equation}
    \begin{aligned}
      (II) 
      =&\int_{0}^{\infty}\frac{e^{t}}{\sqrt{2\pi t}}e^{-\mu x}dx\int_{0}^{\infty}e^{-\frac{[u+(z-x+\lambda t)]^{2}}{2t}}du\\
      =&e^{(1-\frac{\lambda^{2}}{2})t}\int_{0}^{\infty}e^{-(\mu-\lambda) x}dx\int_{0}^{\infty}\frac{1}{\sqrt{2\pi t}}e^{-\frac{[u+(z-x)]^{2}}{2t}}e^{-\lambda u-\lambda z}du.
    \end{aligned}
  \end{equation}
Note that $\int_{0}^{\infty}e^{-(\mu-\lambda) x}dx\int_{0}^{\infty}\frac{1}{\sqrt{2\pi t}}e^{-\frac{[u+(z-x)]^{2}}{2t}}e^{-\lambda u-\lambda z}du$ is finite as $\mu > \lambda$; therefore, one can find a constant $c_{2}$ depending on $z$ such that the above is bounded by $ c_{2}e^{(1-\frac{\lambda^{2}}{2})t}.$ 
By a similar computation used to prove the bound on $(I)$, we get
  \begin{equation}
    \mathbb{P}\biggl(\exists i\in I, k \le {n}^{i}(t):p_{i}+\chi^{i}_{k}(t)-\lambda t \ge z, p_{i}<-(\mu - \lambda)t-\delta\sqrt{t}\biggr) \le c_{3}\int_{\delta}^{\infty}e^{-\frac{u^{2}}{2}}du. \label{2.20}
  \end{equation}
  Therefore
     \begin{align*}
      &\mathbb{P} \biggl(\exists i\in \mathbb N, k \le n^{i}(t):p_{i}+\chi_k^i(t)-\lambda t \ge z,p_i\notin[-(\mu - \lambda)t-\delta\sqrt{t},-(\mu -\lambda)t+\delta\sqrt{t}]\biggr)\\
  \le&c_{1}\int_{-(\mu-\lambda)\sqrt{t}}^{-\delta}\frac{1}{\sqrt{2\pi}}e^{-\frac{u^{2}}{2}}du +c_{2}e^{(1-\frac{\lambda^{2}}{2})t}+ c_{3}\int_{\delta}^{\infty}e^{-\frac{u^{2}}{2}}du,
     \end{align*}
which concludes the proof of this lemma once we choose $\delta, t_0$ sufficiently large.
\end{proof}

\end{document}
